% Generated draft: requires mathematical and source review.
\documentclass[11pt]{article}
\usepackage[margin=25mm]{geometry}
\usepackage{fontspec}
\setmainfont{Latin Modern Roman}
\usepackage{amsmath,amssymb,mathtools,unicode-math}
\setmathfont{Latin Modern Math}
\usepackage{enumitem,needspace,xurl,hyperref}
\setlength{\parindent}{0pt}
\setlength{\parskip}{5pt}
\setlength{\emergencystretch}{3em}
\newcommand{\N}{\mathbb{N}}
\newcommand{\Z}{\mathbb{Z}}
\newcommand{\Q}{\mathbb{Q}}
\newcommand{\R}{\mathbb{R}}
\newcommand{\C}{\mathbb{C}}
\newcommand{\F}{\mathbb{F}}
\newcommand{\A}{\mathbb{A}}
\newcommand{\PP}{\mathbb{P}}
\newcommand{\E}{\mathbb{E}}
\newcommand{\eps}{\varepsilon}
\newcommand{\dd}{\,\mathrm{d}}
\begin{document}
\section*{UnsolvedMath ID: 30001130}
\textbf{Source problem code:} catalogue entry
\subsection{Definitions and mathematical statement}
The record poses two questions about toric fibers and displacement by probes. Its geometric conventions are not fully specified. The following gives an explicit formulation under the standard compact smooth toric interpretation; the choice of this setting, the meaning of probe, and the Fano normalization require source review.

Let \((M^{2n},\omega)\) be a compact connected symplectic manifold with an effective Hamiltonian action of \(T^n=\R^n/\Z^n\), and let \(\mu:M\to\R^n\) be a moment map with image \(P\). Assume that \(P\) is a full-dimensional Delzant polytope: it is simple, and the primitive inward integral normals of the facets meeting at each vertex form a basis of \(\Z^n\). For \(x\in\operatorname{int}P\), write \(L_x=\mu^{-1}(x)\) for the corresponding Lagrangian torus. Define Hamiltonian nondisplaceability by
\[
 \forall\phi\in\operatorname{Ham}(M,\omega),\qquad
 \phi(L_x)\cap L_x\ne\varnothing,
\]
where \(\operatorname{Ham}(M,\omega)\) denotes the group of time-one maps of Hamiltonian isotopies.

For an explicit ordinary-probe interpretation, choose a facet \(F\) with primitive inward integral normal \(\eta_F\), a point \(b\in\operatorname{relint}F\), and a primitive vector \(v\in\Z^n\) with \(\langle\eta_F,v\rangle=1\). Let
\[
 \ell(b,v)=\sup\{t\geq0:b+sv\in P\text{ for every }0\leq s\leq t\}.
\]
The associated probe follows \(b+tv\) for \(0\leq t<\ell(b,v)\). Define its strict first-half region, and the region missed by all such probes, by
\[
 \mathcal D_{\mathrm{pr}}(P)
 =\bigcup_{F,b,v}\{b+tv:0<t<\ell(b,v)/2\},
 \qquad
 \mathcal R_{\mathrm{pr}}(P)
 =\operatorname{int}P\setminus\mathcal D_{\mathrm{pr}}(P).
\]
Here the union ranges over the choices just specified. This definition makes explicit one conventional meaning of displacement by a probe. Failure to belong to \(\mathcal D_{\mathrm{pr}}(P)\) is a failure of that probe criterion; it does not by definition imply Hamiltonian nondisplaceability. The source does not specify whether any enlarged class of probes is intended, or how points at exactly half the probe length are treated.

\begin{enumerate}
\item For \(n=2\), whenever \(U\subseteq\mathcal R_{\mathrm{pr}}(P)\) is a nonempty open subset of \(\operatorname{int}P\), is
\[
 \forall x\in U\ \forall\phi\in\operatorname{Ham}(M,\omega),
 \qquad \phi(L_x)\cap L_x\ne\varnothing?
\]
Can this nondisplaceability be detected by a suitable Floer homology? A possible precise form of the detection question asks, for each such \(x\), for admissible Floer-theoretic data \(\mathcal B_x\) such that a well-defined group
\[
 HF(L_x,L_x;\mathcal B_x)\ne0
\]
provides a Hamiltonian nondisplaceability obstruction. The symbol \(\mathcal B_x\) is a placeholder: the source does not specify coefficients, bulk deformations, bounding cochains, unobstructedness assumptions, or the required detection theorem. Thus this Floer-theoretic target remains unresolved at the level of formulation.

\item For the standard normalized smooth Fano interpretation, suppose that there is a distinguished lattice point \(p_*\in\operatorname{int}P\cap\Z^n\), that it is the unique interior lattice point, and that
\[
 P=\{x\in\R^n:\langle\eta_F,x-p_*\rangle\geq-1
                    \text{ for every facet }F\}.
\]
In arbitrary dimension \(n\), is
\[
 \operatorname{int}P\setminus\{p_*\}
 \subseteq\mathcal D_{\mathrm{pr}}(P)?
\]
This is an interior-fiber formulation of the supplied Fano question. The record says ``every point'' of a Fano moment polytope, without explicitly restricting to its interior. Boundary fibers are not Lagrangian tori, and the probe region defined above lies in the interior. Whether boundary points are intended is therefore an unresolved qualification. The record also does not supply the normalization identifying its special interior integral point, so the displayed Fano normalization is conditional rather than a verified source hypothesis.
\end{enumerate}
\subsection{Short English statement}
For UnsolvedMath ID \(30001130\), the first question asks whether open regions of a two-dimensional moment polytope missed by probes correspond to Hamiltonian nondisplaceable torus fibers, and whether suitable Floer homology detects this. The second asks whether probes displace every other interior fiber of a normalized Fano polytope. The probe convention, geometric hypotheses, Fano normalization, treatment of boundary points, and precise Floer theory require review. The supplied partially solved label is dated curation, not an independently verified conclusion.
\subsection{Sources}
\begin{itemize}
\item Supplied local UnsolvedMath registry record, permanent ID \(30001130\), title: \emph{Probe Displaceability and Floer Detection of Toric Fibers}; source problem code \texttt{catalogue entry}; collection: Oberwolfach Reports Open Problems; proposed year: 2009. No individual proposer is supplied.
\item Supplied authoritative report citation: \emph{Toric Geometry} (2009), Oberwolfach Reports, DOI 10.4171/owr/2009/01, pp.~5--74. Source URL and supplied literature link labeled ``original OWR report'': \url{https://doi.org/10.4171/owr/2009/01}. The background describes Oberwolfach Reports as documenting workshops and their open-problem discussions. No problem-specific page or passage is supplied.
\item Supplied literature link labeled ``McDuff probes'': \url{https://arxiv.org/abs/0904.1686}. No fuller bibliographic citation is supplied.
\item Supplied literature link labeled ``toric Floer bulk deformations'': \url{https://arxiv.org/abs/0810.5654}. No fuller bibliographic citation is supplied.
\item Literature assessment date: \texttt{2026-08-21}. This is supplied dated curation, not independently verified literature assessment. The supplied status is \texttt{partially\_solved}. The assessment states that the special Fano fiber is known to be nondisplaceable, that probe and Floer methods settle broad families and low dimensions, and that universal probe and higher-dimensional nondisplaceability questions remain open. The background calls this a dated literature triage derived from the supplied Oberwolfach curation and explicitly requests independent verification of time-sensitive status and cited claims.
\item The registry marks the statement as \texttt{exact} and reports a manual review finding the extraction self-contained, mathematically coherent, and free of report-layout artifacts. These are supplied verification claims; they do not independently establish the geometric conventions or mathematical status used here.
\end{itemize}
\end{document}
